% 六套模板共用同一正文骨架；换版式时保留本文件及全部实测材料。
\section{引言}
\ReportTODO{第一段：介绍实验背景、研究对象及意义。}

\ReportTODO{第二段：介绍本次实验内容、目的、研究问题、测量对象和实验思路。}

\section{原理}
\subsection{实验原理}
\ReportTODO{归纳实际涉及的物理规律、理论模型、关键公式和适用条件，说明符号、单位及理论量与可测量之间的关系。}
公式须编号并定义符号。例如，模型 $y=f(x_1,\ldots,x_n)$ 在线性近似且输入相互独立时，合成标准不确定度为
\begin{equation}
u_c^2(y)=\sum_{i=1}^{n}
\left(\frac{\partial f}{\partial x_i}\right)^2 u^2(x_i).
\label{eq:uncertainty}
\end{equation}
式\eqref{eq:uncertainty}为示例。\ReportTODO{替换为本实验的实际模型；输入相关时加入协方差项，不适用的示例应删除。}

\subsection{实验装置}
\ReportTODO{列出实际仪器型号、量程、分辨率、连接关系、环境和样品状态；未知信息标注待补充，不得估造。}

\subsection{实验方法}
\subsubsection{步骤1：按实际操作命名}
\ReportTODO{按实际顺序说明装置设置、关键操作、观测对象、记录内容和重复次数。}
\subsubsection{步骤2：按实际操作命名}
\ReportTODO{根据实验的实际步骤增删本级标题；不得虚构或机械照抄未执行的操作。}

\section{结果与分析讨论}
\subsection{定量结果与数据处理}
\ReportTODO{区分原始、计算、理论和估计值，给单位、有效数字及代表性计算。表\ref{tab:raw}仅为格式示例，不含实测数据。}
\begin{table}[H]
\centering
\caption{原始数据记录表}
\label{tab:raw}
\begin{tabularx}{\columnwidth}{@{}cXXX@{}}
\toprule
\ReportTODO{列名及单位} & \ReportTODO{实测值} & \ReportTODO{实测值} & \ReportTODO{备注} \\
\midrule
\ReportTODO{仅填写真实原始读数} & & & \\
\bottomrule
\end{tabularx}
\end{table}
\ReportTODO{说明定标、拟合、统计或插值方法；给出模型、参数、单位及适用的 R\textsuperscript{2}、残差或不确定度。图\ref{fig:result}须由真实数据生成。}
\begin{figure}[H]
\centering
\fbox{\parbox[c][22mm][c]{0.9\columnwidth}{\centering 待补充：根据真实数据生成的图}}
\caption{实验结果图}
\label{fig:result}
\end{figure}

\subsection{规律、理论比较与误差讨论}
\ReportTODO{结合实验原理解释实测规律，与理论预期或标准值比较；分析偏差和误差来源，保留异常值并说明必要的剔除判据。至少提出两条对应误差来源、可执行的改进建议。}

\section{实验总结}
\ReportTODO{回应实验目的，总结有证据支持的主要结果和结论，并说明实验局限及改进方向；不引入未经论证的新结果。}

\renewcommand{\refname}{五、参考文献}
\bibliographystyle{gbt7714-numerical}
\bibliography{report_template/references}
\ifdefined\ReportEndColumns\ReportEndColumns\fi
\ReportEnglishSummary

\ifReportAppendix
\section*{附录：程序与补充说明}
\ReportTODO{仅按需附程序、推导或其他非实验数据补充说明；实验数据另存于本地文件，不列入附录。附录置于英文摘要和英文关键词之后。}
\fi
